% !TeX program = xelatex
% 2026 CUMCM 电子版论文：第一页必须是摘要专用页。
% 编译：xelatex paper_v1.tex -> bibtex paper_v1 -> xelatex x2
% 本文件由 main.tex + sections/*.tex 自动合并生成；修改请改 sections 源文件后用 _merge.py 重新生成
\documentclass[zihao=-4,a4paper,UTF8,fontset=none]{ctexart}

% template
\usepackage{cumcm2026}

% set English and Number font
\setmainfont{TimesNewRoman.ttf}[
  Path=fonts/,
  BoldFont=TimesNewRomanBold.ttf,
  ItalicFont=TimesNewRomanItalic.ttf,
  BoldItalicFont=TimesNewRomanBoldItalic.ttf
]

% set Chinese font
\setCJKmainfont[
  Path=fonts/,
  BoldFont=SimHei.ttf
]{SimSun.ttc}

\setCJKfamilyfont{zhsong}[
  Path=fonts/
]{SimSun.ttc}

\setCJKfamilyfont{zhhei}[
  Path=fonts/,
  AutoFakeBold=2.5
]{SimHei.ttf}

% 代码环境(listings \ttfamily)中的中文使用宋体
\setCJKmonofont[
  Path=fonts/
]{SimSun.ttc}

\providecommand{\songti}{}
\providecommand{\heiti}{}
\renewcommand{\songti}{\CJKfamily{zhsong}}
\renewcommand{\heiti}{\CJKfamily{zhhei}}

%--------------------------------------------------------------------
% 标题与关键词（团队最终确认）
\title{基于储能套利与滚动优化的微网电力调控策略研究}
\newcommand{\keywords}{微网调度；储能套利；线性规划；日前优化；敏感性分析}

\begin{document}

% 2026 电子提交版不得包含承诺书和编号专用页；第一页直接为摘要页。
\maketitle

\begin{cumcmabstract}
% 摘要为全队定稿部分：Q1 内容已按冻结结果写入，Q2--Q4 部分待对应问题求解完成后补充修改。
针对微网与外部电网的电力调控问题，本文围绕“购电量与储能充放电量的最优调度”这一核心任务，构建了从确定性日前优化到不确定性滚动调整的递进式优化模型体系。

针对问题一，在电价、小区负载均已知且储能日初日末电量相等的确定性条件下，建立了以全天购电费用最小为目标、包含电力平衡、储能状态转移、容量与功率限制及边界条件的线性规划模型；利用储能“低价充电、高价放电”的套利机理，通过 HiGHS 求解器求得全局最优解：全天购电量为 $59{,}482.70$ kWh，全天购电费为 $35{,}126.95$ 元。与不使用储能的基线方案（$48{,}052.05$ 元）相比，储能运行带来净节省 $12{,}925.10$ 元，节省比例达 $26.9\%$。逐时段约束验证表明模型解严格满足全部约束，能量守恒误差为 $0.0000$ kWh，验证了模型与求解的正确性。

针对问题二，……（待 Q2 求解后补充）。针对问题三，……（待 Q3 求解后补充）。针对问题四，……（待 Q4 求解后补充）。

\cumcmkeywords{\keywords}
\end{cumcmabstract}

% 问题重述
\section{问题重述}

\subsection{问题背景}

微网是为促进光伏等分布式能源就地消纳、缓解并网压力而发展起来的小型发用电系统。非孤岛运行的微网可以与外部电网进行电力交易：当外网电价较低或分布式电源出力富余时，为储能设备充电；当外网电价较高时，释放储存的电能。微网调控的核心任务是，在小区负载、光伏出力与储能电量等信息的支撑下，给出各时段的购电量与储能充放电量，使微网供电始终满足小区负载需求，同时尽可能降低购电费用。

某小区微网由光伏发电系统、储能设备与小区负载构成。储能设备最大容量为 $12000$ kWh，最大充放电功率为 $5000$ kW，充放电效率为 $90\%$，其电量须保持在 $1200$--$10800$ kWh 之间；2025年1月1日 0:00 的初始电量为 $6000$ kWh。题目提供了小区一天内以 10 分钟为间隔的电价、负载与光伏功率数据（附件1），以及全年的实际运行与预报数据（附件2--4），要求制定不同条件下的购电策略。

\subsection{问题内容}

\textbf{问题一：} 在每天电价和小区负载相同、光伏发电功率预测已知的确定性条件下，要求微网提供的电能不低于小区负载，且储能设备在 0:00 与 24:00 的储电量相同，建立每天 0:00 制定当天计划购电策略的数学模型；基于附件1与附录1的数据，给出指定时间段的购电量、全天购电量与购电费（表1），以及储能设备分时段充放电量与边界储电量（表2），并将完整策略写入 result1.xlsx。

\textbf{问题二：} 在每天电价相同、负载与光伏功率逐日变化的条件下，若微网供电低于负载，需向外网紧急购电（电价为交易时刻电价的 5 倍），建立每天 0:00 的计划购电模型，并基于附件1、附件2给出全年策略及指定日期的结果。

\textbf{问题三：} 在每天 0:00、6:00、12:00 与 18:00 可获取未来 24 小时光伏功率预报的条件下，允许根据新预报调整购电策略（调减部分收取 $50\%$ 违约费用，调增部分按 $1.5$ 倍电价计费），建立含计划、紧急与调整费用的购电模型，并分析引入滚动调整的必要性。

\textbf{问题四：} 在电价实时波动的条件下，重建购电策略模型，并在波动电价下重新求解问题二与问题三。

四个问题层层递进：问题一建立确定性日前优化的基准模型；问题二引入实际运行中的供需偏差与紧急购电机制；问题三进一步引入多时点预报与策略调整；问题四将确定电价推广为波动电价。本文依次建模求解，并在最后检验模型的有效性与稳健性。

% 问题分析
\section{问题分析}

\subsection{数据理解}

附件1 提供了某一天以 10 分钟为间隔的电价、小区负载与光伏发电预测功率数据，共 144 个时段（其中 ``0:00+1'' 表示当天 24:00）。对附件1 进行统计分析，结果如表~\ref{tab:data_stat} 所示。

\begin{table}[H]
  \centering
  \caption{附件1 主要变量的描述性统计}
  \label{tab:data_stat}
  \begin{tabular}{ccccccc}
    \toprule
    变量 & 均值 & 最小值 & 最大值 & 峰谷比/波动 & 单位 \\
    \midrule
    电价 & 0.766 & 0.371 & 1.395 & 3.76 & 元/kWh \\
    小区负载 & 4626.0 & 3309.4 & 5959.0 & 1.80 & kW \\
    光伏发电预测功率 & 2311.8 & 0.0 & 7612.3 & -- & kW \\
    \bottomrule
  \end{tabular}
\end{table}

数据呈现三个明显特征。其一，电价具有显著的分时结构：最低电价 $0.371$ 元/kWh 与最高电价 $1.395$ 元/kWh 相差约 3.76 倍，为储能“低价充电、高价放电”的套利运行提供了经济空间。其二，小区负载日间波动明显（$3309$--$5959$ kW），存在早晚两个用电高峰。其三，光伏发电集中在白天（全天 89 个时段出力为正，占 $61.8\%$），峰值功率 $7612.3$ kW 超过负载峰值，光伏富余时段需要储能吸收多余电量。此外，数据完整无缺失、数值均在合理范围，可直接用于建模。

\subsection{总体思路}

四个子问题共享同一个物理系统与调控目标，区别在于信息条件与不确定性逐问增强，本文据此构建递进式的优化模型体系：

\textbf{问题一}是确定性日前优化：电价、负载与光伏出力均精确已知，且要求储能日初日末电量相等、实现日循环运行。其本质是一个储能套利调度问题，目标函数与全部约束均为决策变量的线性函数，故建立线性规划模型，借助成熟的单纯形类求解器即可获得全局最优解。该模型同时作为后续问题的基准框架。

\textbf{问题二}在问题一的基础上引入负载与光伏的实际波动以及“供电不足时按 5 倍电价紧急购电”的惩罚机制，需在日前计划中考虑不确定性的代价。

\textbf{问题三}引入 0:00、6:00、12:00、18:00 四个预报时刻，将问题转化为含计划、调整与违约费用的滚动优化问题。

\textbf{问题四}将确定电价替换为实时波动电价，检验前述模型在更一般价格信号下的适应性。

四个问题之间的逻辑关系与数据流如图~\ref{fig:roadmap} 所示（技术路线图待 A/B 确认后由 DrawIO 生成并插入）。

\begin{figure}[H]
  \centering
  % TODO: 技术路线图由 4drawio 阶段生成后插入，此处暂以文字占位，不占用最终图号
  \fbox{\parbox{0.8\textwidth}{\centering 技术路线图（待插入）}}
  \caption{总体技术路线}
  \label{fig:roadmap}
\end{figure}

\subsection{问题一的分析}

问题一的决策内容是全天 144 个时段的购电量与储能充放电量。需要同时处理三个相互耦合的环节：\textbf{电力平衡}——任一时刻购电、光伏与储能放电之和须覆盖负载与储能充电需求；\textbf{储能动态}——充放电改变储电量，且充放电过程伴随效率损耗；\textbf{经济目标}——在满足供电的前提下最小化购电费用。储能的存在使各时段决策相互关联：当前时段充入的电量可以在后续高价时段释放，从而将购电需求从高价时段“转移”到低价时段，这正是储能套利的机理。因此问题一可自然表述为一个多时段耦合的线性规划问题：以各时段购电量与充放电量为决策变量，以全天购电费用为目标，以电力平衡、储能状态转移、容量与功率限制、边界条件为约束，求解器给出的即为理论最优的日前调度策略。

依据题意，问题一还需满足若干结构性条件：储能电量运行区间为 $1200$--$10800$ kWh（附录1 硬性要求）；充放电效率 $90\%$ 按字面理解为充电、放电单侧效率各为 $90\%$，即往返效率 $81\%$；0:00 与 24:00 储电量均为初始值 $6000$ kWh；微网不允许向外网售电，购电量非负，光伏富余时允许弃光。

% 模型假设
\section{模型假设}

本文在建模过程中采用以下假设：

\noindent 1) \textbf{确定性信息假设（问题一）}：一天内各时段的电价、小区负载与光伏发电功率均精确已知，且 10 分钟内保持不变。附件1 以 10 分钟为间隔给出数据，该假设与数据粒度一致，是日前优化模型的通用前提。

\noindent 2) \textbf{效率口径假设}：储能充放电效率按题目表述取单侧效率 $\eta = 90\%$，即充电输入 $E_{\mathrm{charge}}$ 时仅 $\eta E_{\mathrm{charge}}$ 存入电池，电池放出 $E_{\mathrm{discharge}}$ 时仅 $\eta E_{\mathrm{discharge}}$ 供给负载，往返效率为 $\eta^2 = 81\%$。效率在问题一的时间尺度内视为恒定，不随充放电功率与荷电状态变化。

\noindent 3) \textbf{储能物理约束假设}：储能设备同一时段不同时充电与放电。由于购电价格为正面充放电存在损耗，同时充放电只会产生净损耗而不降低购电费用，最优解将自然规避该情形，故模型中无需显式加入互斥约束（后文给出严格论证）。

\noindent 4) \textbf{电量平衡与边界假设}：微网供电允许大于负载需求，即光伏出力超过“负载＋储能可吸收电量”时允许弃光，弃光不产生收益也不产生成本；储能设备在 0:00 与 24:00 的储电量均为 $6000$ kWh，实现日循环运行。

\noindent 5) \textbf{交易规则假设}：微网只能从外网购电，不向外网售电；购电费用等于购电量与对应时段电价的乘积，无其他附加费用；问题一暂不考虑储能设备折旧、电网功率约束与电能质量要求等次要因素。

\noindent 6) \textbf{数据有效性假设}：赛题提供的附件数据真实有效，单位统一，无缺失与异常值（经数据审查验证）。

% 符号说明
\section{符号说明}

本文用到的主要符号及其含义如表~\ref{tab:symbols} 所示。符号体系统一在全部子问题中沿用，仅问题三、问题四增加少量带上下标的变量（届时在对应章节补充说明）。

\begin{table}[H]
  \centering
  \caption{主要符号与定义}
  \label{tab:symbols}
  \begin{tabularx}{0.92\textwidth}{C{0.22\textwidth}Y C{0.14\textwidth}}
    \toprule
    符号 & 含义 & 单位 \\
    \midrule
    $t$ & 时段索引，$t=1,2,\ldots,144$，对应 10 分钟间隔 & -- \\
    $\Delta t$ & 时间间隔长度，$\Delta t=1/6$ & h \\
    $Price(t)$ & 时段 $t$ 的电价 & 元/kWh \\
    $Load(t)$ & 时段 $t$ 的小区负载功率 & kW \\
    $PV(t)$ & 时段 $t$ 的光伏发电预测功率 & kW \\
    $E_{buy}(t)$ & 时段 $t$ 从外网购买的电量 & kWh \\
    $E_{charge}(t)$ & 时段 $t$ 的名义充电能量（输入侧，入池 $\eta E_{charge}$） & kWh \\
    $E_{discharge}(t)$ & 时段 $t$ 的名义放电能量（电池侧，供负载 $\eta E_{discharge}$） & kWh \\
    $E_{storage}(t)$ & 时刻 $t$ 结束时储能设备的储电量 & kWh \\
    $E_{\min}$, $E_{\max}$ & 储能电量下限、上限（$1200$, $10800$） & kWh \\
    $P_{\max}$ & 储能最大充放电功率（$5000$） & kW \\
    $\eta$ & 储能单侧充放电效率（$0.9$，往返 $0.81$） & -- \\
    $Z$ & 全天购电费用（目标函数值） & 元 \\
    \bottomrule
  \end{tabularx}
\end{table}

\noindent 说明：时段内变化的量（购电量、充放电量）以下标 $t$ 对应时段编号；时刻定义的量（储电量）对应时间点编号，$E_{storage}(0)$ 与 $E_{storage}(144)$ 分别表示 0:00 与 24:00 的储电量。功率与电量的换算关系为电量 $=$ 功率 $\times \Delta t$。

% 问题一：确定性日前购电优化模型
\section{问题一：确定性条件下的日前购电优化模型}

\subsection{建模思路}

问题一在电价、小区负载与光伏出力全部精确已知的条件下，求解全天 144 个时段的购电量与储能充放电量，使购电费用最小。这是典型的微网确定性经济调度问题，其决策结构与光伏--电池混合系统的需求侧管理问题一致\cite{wu2015demand}：储能设备扮演能量“时间搬运”的角色，在低价时段吸收外网电量与光伏富余电量，在高价时段释放，从而将购电需求从高价时段转移到低价时段。

该问题的目标函数（购电费用为购电量与电价的乘积之和）与全部约束（电力平衡、储能状态转移、容量与功率限制）均为决策变量的线性函数，且决策变量为连续变量，因此可直接建立线性规划（LP）模型。微网能量管理领域的综述研究表明，对于确定性、单目标的经济调度问题，线性规划凭借其全局最优性和求解效率是公认的首选方法\cite{zia2018microgrids}，无需引入元启发式算法。问题一的最优解即为微网在确定性信息下可达到的理论最优日前策略，同时也是问题二至问题四建模的基准框架。

\subsection{数据预处理}

附件1 提供了一天 144 个时段的电价、小区负载与光伏发电预测功率，时间间隔 $\Delta t = 10\ \mathrm{min} = 1/6$ h，表中 ``0:00+1'' 即当天 24:00。预处理步骤如下：

\textbf{(1) 时间索引化。} 将 144 个时段依次编号 $t = 1,2,\ldots,144$，其中 $t=1$ 对应 0:00--0:10，$t=61$ 对应 10:00--10:10，$t=144$ 对应 23:50--24:00。

\textbf{(2) 数据完整性校验。} 检查确认附件1 共 144 行、字段齐全，电价、负载与光伏功率均无缺失、无穷或非法数值，电价严格为正，负载与光伏功率非负，时间序列自 0:00 起严格等间隔 10 分钟且顺序正确。

\textbf{(3) 单位统一。} 模型以电量为计量单位，将功率数据乘以时间间隔换算为电量：$Load(t)\cdot\Delta t$、$PV(t)\cdot\Delta t$ 分别表示时段 $t$ 的负载需求电量与光伏发电电量（kWh）。储能每时段最大充放电电量为 $P_{\max}\Delta t = 5000/6 = 833.33$ kWh。

\subsection{储能充放电的线性规划模型}

\subsubsection{决策变量}

对每个时段 $t\in\{1,2,\ldots,144\}$ 定义四类连续决策变量：

\begin{itemize}
  \item $E_{buy}(t)$：时段 $t$ 从外网购买的电量（kWh）；
  \item $E_{charge}(t)$：时段 $t$ 储能的名义充电能量（kWh），即充电输入侧能量，经效率 $\eta$ 折算后实际入池 $\eta E_{charge}(t)$；
  \item $E_{discharge}(t)$：时段 $t$ 储能的名义放电能量（kWh），即电池放出侧能量，经效率 $\eta$ 折算后实际供负载 $\eta E_{discharge}(t)$；
  \item $E_{storage}(t)$：时刻 $t$ 结束时储能的储电量（kWh），$t=1,\ldots,144$。
\end{itemize}

共 $144\times 4 = 576$ 个决策变量；$E_{storage}(0)=6000$ kWh 为已知参数。效率 $\eta=0.9$ 统一作用于能量流出方向：充电时输入 $E_{charge}$ 损失 $10\%$、入池 $\eta E_{charge}$；放电时电池放出 $E_{discharge}$ 损失 $10\%$、供负载 $\eta E_{discharge}$，往返效率为 $\eta^2=0.81$。

\subsubsection{目标函数}

以全天购电费用最小为目标：

\begin{equation}
  \min Z = \sum_{t=1}^{144} Price(t)\cdot E_{buy}(t)
  \label{eq:obj}
\end{equation}

光伏发电与储能放电均不产生购电成本，模型将优先利用这两类零成本能量。

\subsubsection{约束条件}

\textbf{(1) 电力平衡约束。} 任一时刻微网供给能力须不低于负载需求：

\begin{equation}
  E_{buy}(t) + PV(t)\Delta t + \eta E_{discharge}(t) \ge Load(t)\Delta t + E_{charge}(t),\quad t=1,\ldots,144
  \label{eq:balance}
\end{equation}

其中左侧为供给（购电、光伏发电、储能放电），右侧为需求（负载用电、储能充电输入）。约束取不等式，允许光伏出力超过“负载与储能吸收能力之和”时弃光；由于目标只惩罚购电量且无售电收益，最优解不会故意产生多余供给。

\textbf{(2) 储能状态转移方程。} 相邻时刻储电量满足：

\begin{equation}
  E_{storage}(t) = E_{storage}(t-1) + \eta E_{charge}(t) - E_{discharge}(t),\quad t=1,\ldots,144
  \label{eq:state}
\end{equation}

\textbf{(3) 容量约束。} 储电量须始终处于题目规定的安全运行区间：

\begin{equation}
  1200 = E_{\min} \le E_{storage}(t) \le E_{\max} = 10800,\quad t=1,\ldots,144
  \label{eq:capacity}
\end{equation}

\textbf{(4) 充放电功率约束。} 名义充放电能量受最大充放电功率限制：

\begin{equation}
  \begin{split}
    0 \le E_{charge}(t) \le P_{\max}\Delta t = 833.33,\quad t=1,\ldots,144,\\
    0 \le E_{discharge}(t) \le P_{\max}\Delta t = 833.33,\quad t=1,\ldots,144
  \end{split}
  \label{eq:power}
\end{equation}

\textbf{(5) 边界条件。} 储能日初、日末电量相同，实现日循环运行：

\begin{equation}
  E_{storage}(0) = 6000,\qquad E_{storage}(144) = 6000
  \label{eq:boundary}
\end{equation}

\textbf{(6) 非负约束。} 微网不向外网售电：

\begin{equation}
  E_{buy}(t) \ge 0,\quad t=1,\ldots,144
  \label{eq:nonneg}
\end{equation}

\subsubsection{充放电互斥性的理论论证}

模型式\eqref{eq:balance}--\eqref{eq:nonneg}为纯线性规划，未显式禁止同时充放电。以下证明最优解不会出现同时充放电，故无需引入 0-1 互斥变量。设某可行解在时段 $t$ 同时有 $E_{charge}(t)=c>0$、$E_{discharge}(t)=d>0$，令 $\delta=\min(c,d)>0$，构造修正解 $\tilde c=c-\delta$、$\tilde d=d-\delta$（其余变量不变）。修正后：电力平衡左侧减少 $\eta\delta$、右侧减少 $\delta$，由 $\eta<1$ 知约束\eqref{eq:balance}的松弛量增加，仍然成立；状态方程\eqref{eq:state}右端增加 $(1-\eta)\delta>0$，储电量增大，容量约束\eqref{eq:capacity}不会因此被破坏；目标函数值不变。故任一含同时充放电的可行解均可经上述修正消去互斥量而不劣化目标，最优解中必存在无同时充放电的调度策略。这一结论也由求解后逐时段诊断得到验证（第~\ref{subsec:verify}节：同时充放电时段数为 0）。

\subsubsection{模型特性}

式\eqref{eq:obj}--\eqref{eq:nonneg}构成一个线性规划问题：576 个连续变量、144 个电力平衡约束、144 个状态转移方程、288 个容量上下界、576 个充放电功率界限、144 个购电非负约束与 2 个边界条件。线性规划是凸优化问题，其局部最优解即全局最优解，无需迭代调参，求解结果可复现、可验证。

\subsection{模型求解}

模型采用 Python 语言 PuLP 建模库构建，调用开源求解器 HiGHS 求解。求解状态为 Optimal，目标值有限，求解耗时约 $0.01$ 秒，得到全局最优的日前购电策略。所得策略满足全部约束（第~\ref{subsec:verify}节给出逐项验证），完整 144 时段的购电量与分时段充放电量已按附件5 模板写入 result1.xlsx。

\subsection{购电与充放电计划}

表~\ref{tab:q1_table1}给出题目要求的 6 个指定时段的购电量及全天汇总，表~\ref{tab:q1_table2}给出 6 个四小时时段的充放电量与边界储电量。

\begin{table}[H]
  \centering
  \caption{微网在指定时间段的购电量及全天的购电量和购电费}
  \label{tab:q1_table1}
  \begin{tabular}{ccccc}
    \toprule
    时间段 & 购电量(kWh) & 时间段 & 购电量(kWh) \\
    \midrule
    10:00--10:10 & 0.00 & 12:00--12:10 & 480.41 \\
    14:00--14:10 & 0.00 & 16:00--16:10 & 445.43 \\
    18:00--18:10 & 531.89 & 20:00--20:10 & 0.00 \\
    \midrule
    \multicolumn{2}{c}{全天购电量(kWh)} & \multicolumn{2}{c}{全天购电费(元)} \\
    \multicolumn{2}{c}{59,482.70} & \multicolumn{2}{c}{35,126.95} \\
    \bottomrule
  \end{tabular}
\end{table}

\begin{table}[H]
  \centering
  \caption{储能设备在指定时间段的充放电量及 0:00 和 24:00 的储电量}
  \label{tab:q1_table2}
  \small
  \begin{tabular}{cccccc}
    \toprule
    时间段 & 充电量(kWh) & 放电量(kWh) & 时间段 & 充电量(kWh) & 放电量(kWh) \\
    \midrule
    0:00--4:00 & 4,500.00 & 0.00 & 4:00--8:00 & 833.33 & 7,073.16 \\
    8:00--12:00 & 4,787.96 & 1,892.22 & 12:00--16:00 & 5,286.04 & 101.22 \\
    16:00--20:00 & 0.00 & 6,422.37 & 20:00--24:00 & 5,333.33 & 3,177.63 \\
    \midrule
    \multicolumn{3}{c}{0:00 储电量(kWh)} & \multicolumn{3}{c}{24:00 储电量(kWh)} \\
    \multicolumn{3}{c}{6,000.00} & \multicolumn{3}{c}{6,000.00} \\
    \bottomrule
  \end{tabular}
\end{table}

表~\ref{tab:q1_table1}显示，10:00--10:10、14:00--14:10 与 20:00--20:10 三个时段购电量为零，供电完全由光伏与储能放电承担；其余三个时段购电量分别为 480.41、445.43 与 531.89 kWh。全天累计购电 $59{,}482.70$ kWh，购电费用 $35{,}126.95$ 元。表~\ref{tab:q1_table2}显示，储能全天充电集中在 0:00--4:00、8:00--16:00 与 20:00--24:00 等电价相对较低的时段，放电集中在 4:00--8:00 与 16:00--20:00 两个用电高峰时段；0:00 与 24:00 储电量均为 $6{,}000.00$ kWh，满足题目要求的边界条件。

\subsection{储能调度策略分析}

图~\ref{fig:soc}给出了最优策略下储能电量随时间的演化轨迹，并与电价曲线对照。可以看到，储能运行呈现出清晰的“低价充电、高价放电”规律：0:00--4:00 夜间低电价时段储电量快速上升；4:00--8:00 早高峰伴随集中放电明显回落，接近运行下限；白天时段储电量围绕高荷电水平波动，多次触及上限 $10800$ kWh；16:00--20:00 晚高峰再次集中放电；20:00--24:00 恢复充电并于 24:00 回到初始值 $6000$ kWh。全天储电量共 21 个时段触及上限、8 个时段触及下限，表明储能容量在峰谷时段得到较充分利用，但未出现越限。

\begin{figure}[H]
  \centering
  \includegraphics[width=0.88\textwidth]{figures/fig2_soc_curve.png}
  \caption{最优策略下的储能电量轨迹与电价曲线}
  \label{fig:soc}
\end{figure}

图~\ref{fig:interval}按 6 个四小时时段汇总了充放电量，与表~\ref{tab:q1_table2}对应。充电高峰出现在 20:00--24:00（$5{,}333.33$ kWh）与 12:00--16:00（$5{,}286.04$ kWh），放电高峰出现在 4:00--8:00（$7{,}073.16$ kWh）与 16:00--20:00（$6{,}422.37$ kWh），充放电时段与电价峰谷结构在时间上高度互补。

\begin{figure}[H]
  \centering
  \includegraphics[width=0.85\textwidth]{figures/fig5_interval_summary.png}
  \caption{分时段充放电量汇总}
  \label{fig:interval}
\end{figure}

\subsection{储能套利效益分析}

为量化储能的贡献，构造不使用储能的基线策略：微网不配置储能时，各时段购电量即为扣除光伏后的净负载需求，其全天购电费用为

\begin{equation}
  Z_{\mathrm{base}} = \sum_{t=1}^{144} Price(t)\cdot\max\!\big(0,\, Load(t)\Delta t - PV(t)\Delta t \big)
  \label{eq:baseline}
\end{equation}

代入附件1 数据计算得基线成本为 $48{,}052.05$ 元。有储能的最优策略购电费用为 $35{,}126.95$ 元，净节省 $12{,}925.10$ 元，节省比例达 $26.9\%$（图~\ref{fig:cost}）。储能带来的效益来自两方面：一是将低价时段电量搬运至高价时段使用（套利收益），二是将白天光伏富余电量储存供高峰使用、减少高价购电（消纳收益）。

\begin{figure}[H]
  \centering
  \includegraphics[width=0.88\textwidth]{figures/fig6_cost_analysis.png}
  \caption{有/无储能购电成本对比与成本结构分析}
  \label{fig:cost}
\end{figure}

从购电成本结构看（表~\ref{tab:cost_struct}），电价 $0.5$--$1.0$ 元/kWh 的中价时段承担了 $78.8\%$ 的购电量与 $86.5\%$ 的购电费用，是成本的主要来源；电价不低于 $1.0$ 元/kWh 的高价时段购电量仅占 $16.2\%$、成本占比 $10.0\%$，说明储能有效压制了高价时段的购电需求；低价时段的购电仅占 $5.0\%$，对应储能集中充电。购电量占比与成本占比的错位直观反映了套利调度的效果：购电费用向低价时段倾斜，高价时段需求被储能放电替代。

\begin{table}[H]
  \centering
  \caption{购电成本结构}
  \label{tab:cost_struct}
  \begin{tabular}{ccccc}
    \toprule
    电价区间 & 购电量(kWh) & 购电量占比 & 成本(元) & 成本占比 \\
    \midrule
    $<0.5$ 元/kWh & 2,968.05 & 5.0\% & 1,241.86 & 3.5\% \\
    $0.5$--$1.0$ 元/kWh & 46,893.86 & 78.8\% & 30,395.67 & 86.5\% \\
    $\ge 1.0$ 元/kWh & 9,620.78 & 16.2\% & 3,489.43 & 10.0\% \\
    \midrule
    合计 & 59,482.70 & 100\% & 35,126.95 & 100\% \\
    \bottomrule
  \end{tabular}
\end{table}

储能运行强度方面，全天名义总充电量为 $20{,}740.67$ kWh，总放电量为 $18{,}666.60$ kWh，充电损耗 $2{,}074.07$ kWh、放电损耗 $1{,}866.66$ kWh，总转换损耗 $3{,}940.73$ kWh；按可用容量 $9600$ kWh 折算的等效充电循环为 $2.16$ 次/天，处于合理运行强度范围。

\subsection{模型验证}
\label{subsec:verify}

为确认求解结果的正确性，对最优解执行分级验证。

\textbf{输入级验证}：附件1 数据 144 行完整，字段齐全，时间序列自 0:00 起严格每 10 分钟一点且顺序正确，无缺失、无穷或非法数值。

\textbf{求解级验证}：HiGHS 返回 Optimal 状态，目标函数值 $35{,}126.95$ 元有限。

\textbf{约束级验证}：对 144 个时段逐项复核电力平衡、状态转移、容量上下界、充放电功率界限、购电非负与终止电量边界，全部硬约束通过，失败项数为 0。

\textbf{能量守恒验证}：按“输入能量 $=$ 输出能量”核对全天能量收支，结果如表~\ref{tab:energy}所示，误差为 $0.0000$ kWh。

\begin{table}[H]
  \centering
  \caption{全天能量平衡验证}
  \label{tab:energy}
  \begin{tabular}{lcc}
    \toprule
    \multicolumn{2}{c}{输入能量(kWh)} & \multicolumn{1}{c}{输出能量(kWh)} \\
    \cmidrule(r){1-2}\cmidrule(l){3-3}
    购电总量 & 59,482.70 & -- \\
    光伏发电总量 & 55,482.84 & -- \\
    初始储能电量 & 6,000.00 & -- \\
    \midrule
    输入合计 & 120,965.53 & -- \\
    \midrule
    -- & -- & 负载用电总量：111,024.81 \\
    -- & -- & 终止储能电量：6,000.00 \\
    -- & -- & 充电损耗：2,074.07 \\
    -- & -- & 放电损耗：1,866.66 \\
    \midrule
    -- & -- & 输出合计：120,965.53 \\
    \bottomrule
  \end{tabular}
\end{table}

\textbf{诊断级检查}：同时充放电时段数为 0，与理论论证一致；全天供需富余量（对应弃光）仅 $0.0022$ kWh，表明光伏与储能调度高度贴合需求，几乎不存在能量浪费。

综上，问题一的最优策略在满足全部约束的前提下实现购电费用最小，求解结果正确可靠。

% 问题二：不确定性条件下的日前购电优化模型（待 Q2 求解完成后由 C 撰写）
\section{问题二：考虑紧急购电的日前购电优化模型}

% TODO(C)：Q2 求解完成后按以下框架撰写：
% 1) 问题二分析：实际负载/光伏与日前计划的偏差，紧急购电惩罚机制（5 倍电价）
% 2) 模型建立：在问题一 LP 基础上引入供需偏差与紧急购电量 E_emergency(t)
% 3) 求解与验证
% 4) 指定日期（3.20 / 6.21 / 9.23 / 12.21）结果表（表1、表2、表3 格式）
% 5) 结果分析
（本节内容待问题二求解结果交接后撰写。）

% 问题三：多时点预报与滚动调整购电模型（待 Q3 求解完成后由 C 撰写）
\section{问题三：多时点预报下的滚动调整购电模型}

% TODO(C)：Q3 求解完成后按以下框架撰写：
% 1) 问题三分析：四个预报时刻、调整购电的违约（50%）与超额（1.5 倍）计费
% 2) 模型建立：计划-调整-紧急三阶段费用结构，滚动优化
% 3) 求解与验证
% 4) 指定日期结果 + 调整策略必要性分析
（本节内容待问题三求解结果交接后撰写。）

% 问题四：波动电价下的购电策略（待 Q4 求解完成后由 C 撰写）
\section{问题四：波动电价下的购电策略模型}

% TODO(C)：Q4 求解完成后按以下框架撰写：
% 1) 问题四分析：波动电价对 Q2/Q3 模型的输入替换与模型适应性
% 2) 波动电价下 Q2 与 Q3 的重新求解
% 3) result4-2.xlsx / result4-3.xlsx 结果分析与对比
（本节内容待问题四求解结果交接后撰写。）

% 模型检验与敏感性分析
\section{模型检验与敏感性分析}

\subsection{问题一模型的约束验证小结}

问题一的求解结果已通过四级验证（第~\ref{subsec:verify}节）：输入数据完整合规；求解器返回 Optimal；144 个时段的电力平衡、储能状态转移、容量与功率约束逐项复核通过；全天能量守恒误差为 $0.0000$ kWh。此外，诊断检查显示无同时充放电时段、供需富余量仅 $0.0022$ kWh，表明最优策略在物理与经济意义上均合理。

\subsection{问题一模型的参数敏感性}

% TODO(C/B)：待 B 交付 sensitivity.md 的 η=0.9 与 η=√0.9 对比结果后，按以下框架撰写：
% 1) 效率解释差异：单侧 90%（往返 81%）与往返 90%（单侧 √0.9≈0.9487）两种口径
% 2) 两种口径下总购电费用、总购电量与储能调度差异
% 3) 结论：主方案更符合题目字面表述且更保守
（本节数据待 B 的敏感性分析结果交接后填写。）

\subsection{模型稳健性讨论}

% TODO(C)：Q2--Q4 完成后，汇总各问的敏感性分析与稳健性检验。
（本节内容待全部子问题求解完成后撰写。）

% 模型评价与推广
\section{模型评价与推广}

\subsection{模型优点}

（本节内容待全部子问题完成后由 C 撰写：线性规划模型的全局最优性、可复现性与求解效率；递进式模型体系对不确定性逐层扩展的适应性；严格的逐时段验证流程等。）

\subsection{模型局限}

（本节内容待全部子问题完成后由 C 撰写：未计入储能循环寿命与折旧成本\cite{reniers2018improving}；未考虑电网功率约束与电能质量要求；确定性假设在真实运行中的偏差等。）

\subsection{模型推广}

（本节内容待全部子问题完成后由 C 撰写：向多微网协同、需求响应、碳交易成本等场景的推广方向。）


% 2026 强制要求：AI 工具使用声明必须置于参考文献之前。
% 措辞由全队核对后定稿；支撑材料须含 AI工具使用详情.pdf。
\CumcmAIUsed{数学模型的辅助审查与验证、程序编写与调试、论文初稿撰写与语言润色}

\nocite{zia2018microgrids,wu2015demand,parra2016effect,hou2023multi}
\bibliographystyle{gbt7714-numerical}
\bibliography{references}

\newpage
\appendix
% 附录一级标题格式(仅针对附录更改，不要改sty文件)
\titleformat{\section}
  {\centering\heiti\zihao{3}\bfseries}
  {附录~\thesection}
  {0.8em}
  {}

% 附录
\section{支撑材料文件列表}

支撑材料已单独压缩为 RAR 或 ZIP 文件，文件清单如表~\ref{tab:support}所示；提交前必须与实际压缩包逐项核对一致。

\begin{longtable}{@{}p{0.08\textwidth}p{0.42\textwidth}p{0.42\textwidth}@{}}
  \caption{支撑材料文件列表}
  \label{tab:support} \\
  \toprule
  序号 & 文件路径 & 内容说明 \\
  \midrule
  1 & \texttt{code/run\_q1.py} & 问题一主程序（调用 step1--step3，输出 result1.xlsx） \\
  2 & \texttt{code/step1\_data.py} & 问题一数据读取与输入验证 \\
  3 & \texttt{code/step2\_solve.py} & 问题一线性规划模型构建与 HiGHS 求解 \\
  4 & \texttt{code/step3\_deliverables.py} & 问题一结果提取、硬约束验证与交接文件生成 \\
  5 & \texttt{code/generate\_figures.py} & 问题一论文图表生成脚本 \\
  6 & \texttt{AI工具使用详情.pdf} & AI 工具使用情况说明（按 2026 年规定） \\
  \bottomrule
\end{longtable}

% TODO(C)：Q2--Q4 的代码与支撑文件随对应问题求解完成逐步补充。

\section{问题一源程序}

建模所用全部源程序完整、可运行，与论文结果一致。

\lstinputlisting[language=Python,caption={问题一主程序 run\_q1.py}]{code/run_q1.py}
\lstinputlisting[language=Python,caption={数据读取与输入验证 step1\_data.py}]{code/step1_data.py}
\lstinputlisting[language=Python,caption={LP 模型构建与求解 step2\_solve.py}]{code/step2_solve.py}
\lstinputlisting[language=Python,caption={结果提取、验证与交接 step3\_deliverables.py}]{code/step3_deliverables.py}
\lstinputlisting[language=Python,caption={论文图表生成 generate\_figures.py}]{code/generate_figures.py}


\end{document}
